\documentclass{beamer}
\usepackage[utf8]{inputenc}
\usepackage{amsmath,amssymb}
\usepackage{bm}
\usetheme{Madrid}

\title{GMRES: procedura di Arnoldi}
\author{}
\institute{}
\date{}

\begin{document}

\frame{\titlepage}


\begin{frame}{Impostazione del problema}
\textbf{Spazio di Krylov di ordine $k$:}
\[
\mathcal{K}_k(A,b)=\mathrm{span}\{\,b,\,Ab,\,A^2 b,\,\dots,\,A^{k-1}b\,\}.
\]
\medskip

\textbf{Obiettivo:} costruire una \emph{base ortonormale}
\[
V_k=\bigl[v_1,\dots,v_k\bigr], \qquad v_i^\top v_j=\delta_{ij},
\]
per lo spazio $\mathcal{K}_k(A,b)$.
\medskip

\textbf{Perché serve:}
\begin{itemize}
  \item stabilità numerica (evita dipendenze quasi-lineari);
  \item proiezioni semplici ($V_k^\top V_k=I$: $V_k$ trasposto $\times V_k=I$);
  \item riduzione di problemi grandi a problemi piccoli tramite proiezione.
\end{itemize}
\end{frame}

% -------------------------------------------------------
\begin{frame}{Algoritmo di Arnoldi (caso generale)}
\textbf{Input:} matrice $A\in\mathbb{R}^{n\times n}$, vettore $b\neq 0$, passi $k$.\\
\textbf{Output:} base ortonormale $V_{k+1}=[v_1,\dots,v_{k+1}]$ e matrice di Hessenberg $\underline H_k$.
\medskip

\textbf{Inizializzazione:} \quad $v_1 = b/\|b\|$.

\medskip
\textbf{Per $j=1,2,\dots,k$:}
\begin{enumerate}
  \item $w \leftarrow A v_j$
  \item \textit{Ortogon. (Gram–Schmidt modificato):}\\
    per $i=1,\dots,j$: \; $h_{ij}=v_i^\top w$, \quad $w \leftarrow w - h_{ij} v_i$
  \item \textit{Normalizzazione:} \quad $h_{j+1,j}=\|w\|$, \; $v_{j+1} = w / h_{j+1,j}$
\end{enumerate}

\medskip
\textbf{Relazione di Arnoldi:}\quad
\[
A V_k = V_{k+1}\,\underline H_k,
\]
dove $\underline H_k\in\mathbb{R}^{(k+1)\times k}$ è \emph{Hessenberg superiore}.
\end{frame}

\begin{frame}{Definizione: matrice di Hessenberg}
\[
H = [h_{ij}] \in \mathbb{R}^{n \times n} \text{ si dice \emph{matrice di Hessenberg superiore}} 
\]
\[
\text{ se } h_{ij} = 0 \text{ per ogni } i > j + 1.
\]

In altre parole, tutti gli elementi al di sotto della prima sotto-diagonale sono nulli.

\[
H =
\begin{bmatrix}
h_{11} & h_{12} & h_{13} & h_{14} \\
h_{21} & h_{22} & h_{23} & h_{24} \\
0      & h_{32} & h_{33} & h_{34} \\
0      & 0      & h_{43} & h_{44}
\end{bmatrix}
\]

\[
\text{Analogamente, } H \text{ si dice \emph{Hessenberg inferiore} se } h_{ij} = 0 \text{ per ogni } i < j - 1.
\]


\end{frame}

% -------------------------------------------------------
\begin{frame}{Proprietà chiave di Arnoldi}
\begin{itemize}
  \item $V_k=[v_1,\dots,v_k]$ è una base ortonormale di $\mathcal{K}_k(A,b)$.
  \item Ogni nuovo passo \emph{estende} la base: i primi $k$ vettori di $V_{k+1}$ coincidono con quelli di $V_k$.
  \item La struttura di Hessenberg di $\underline H_k$ deriva dall'ortogonalizzazione sequenziale.
  \item Facoltativo: una \emph{ri-ortogonalizzazione} (secondo passaggio MGS) migliora la stabilità numerica.
\end{itemize}

\medskip
\textbf{Interpretazione:} $v_{j+1}$ è la componente di $A v_j$ ortogonale a $\mathrm{span}\{v_1,\dots,v_j\}$, poi normalizzata.
\end{frame}

% -------------------------------------------------------
\begin{frame}{Caso simmetrico: iterazione di Lanczos}
Se $A$ è simmetrica (o hermitiana), Arnoldi si \emph{semplifica} a \textbf{Lanczos}:
\begin{align*}
\beta_1 &= \|b\|, \qquad v_1 = b/\beta_1,\\
w &= A v_1,\qquad \alpha_1 = v_1^\top w,\\
w &\leftarrow w - \alpha_1 v_1,\qquad \beta_2=\|w\|,\qquad v_2=w/\beta_2,
\end{align*}
e per $j\ge 2$ si ha la ricorrenza \emph{a tre termini}:
\[
w \leftarrow A v_j - \beta_j v_{j-1},\qquad
\alpha_j = v_j^\top w,\qquad
w \leftarrow w - \alpha_j v_j,\qquad
\beta_{j+1}=\|w\|,\quad v_{j+1}=w/\beta_{j+1}.
\]

\medskip
La matrice proiettata $T_k = V_k^\top A V_k$ è \textbf{tridiagonale} (con diagonale $\alpha_j$ e sotto/sovra-diagonali $\beta_j$).
\end{frame}

% -------------------------------------------------------
\begin{frame}{Quadro geometrico}
\begin{itemize}
  \item $A v_j$ in generale esce da $\mathrm{span}\{v_1,\dots,v_j\}$.
  \item Sottraendo le proiezioni su $v_1,\dots,v_j$ si ottiene un vettore \emph{ortogonale} alla base corrente.
  \item La normalizzazione produce $v_{j+1}$ a norma unitaria.
  \item Dopo $k$ passi, $V_k$ è una base ortonormale di $\mathcal{K}_k(A,b)$.
\end{itemize}
\end{frame}

% -------------------------------------------------------
\begin{frame}{Riepilogo}
\begin{block}{Definizione}
$\mathcal{K}_k(A,b)=\mathrm{span}\{b,Ab,\dots,A^{k-1}b\}$.
\end{block}

\begin{block}{Arnoldi (generale)}
Costruisce $V_{k+1}$ ortonormale e $\underline H_k$ di Hessenberg tale che $A V_k = V_{k+1}\underline H_k$.
\end{block}

\begin{block}{Lanczos (simmetrico)}
Ricorrenza a tre termini; matrice proiettata tridiagonale $T_k$.
\end{block}

\begin{block}{Note pratiche}
Usare Gram–Schmidt modificato (con possibile ri-ortogonalizzazione) per robustezza numerica.
\end{block}
\end{frame}


\begin{frame}{Esempio: Processo di Arnoldi su una matrice $3\times3$}
\small

Consideriamo
\[
A=\begin{bmatrix}
2 & 1 & 0\\
1 & 2 & 1\\
0 & 1 & 2
\end{bmatrix},
\qquad
b=\begin{bmatrix}1\\0\\0\end{bmatrix}.
\]

Vogliamo costruire una base ortonormale $\{v_1,v_2,v_3\}$ per lo spazio di Krylov
\[
\mathcal{K}_3(A,b)=\text{span}\{b,Ab,A^2b\}.
\]

\textbf{Inizializzazione:}
\[
v_1 = \frac{b}{\|b\|} =
\begin{bmatrix}1\\0\\0\end{bmatrix}.
\]
\end{frame}


\begin{frame}{Iterazione $j=1$:}
\[
w = A v_1 =
\begin{bmatrix}2\\1\\0\end{bmatrix},
\quad h_{11}=v_1^\top w=2,
\]
\[
w \leftarrow w - h_{11}v_1 = \begin{bmatrix}0\\1\\0\end{bmatrix},
\quad h_{21}=\|w\|=1,
\quad v_2 = \frac{w}{h_{21}}=\begin{bmatrix}0\\1\\0\end{bmatrix}.
\]

\end{frame}
\begin{frame}{Iterazione $j=2$:}
\[
w = A v_2 =
\begin{bmatrix}1\\2\\1\end{bmatrix},
\]
\[
h_{12}=v_1^\top w=1 \;\Rightarrow\; w \leftarrow w-h_{12}v_1=\begin{bmatrix}0\\2\\1\end{bmatrix},
\]
\[
h_{22}=v_2^\top w=2 \;\Rightarrow\; w \leftarrow w-h_{22}v_2=\begin{bmatrix}0\\0\\1\end{bmatrix},
\]
\[
h_{32}=\|w\|=1,\quad v_3=\frac{w}{h_{32}}=\begin{bmatrix}0\\0\\1\end{bmatrix}.
\]

\end{frame}
\begin{frame}{Iterazione $j=3$:}
\[
w = A v_3 =
\begin{bmatrix}0\\1\\2\end{bmatrix},
\quad h_{23}=v_2^\top w=1,
\quad w \leftarrow w - h_{23}v_2 = \begin{bmatrix}0\\0\\2\end{bmatrix},
\]
\[
h_{33}=v_3^\top w=2,
\quad w \leftarrow w - h_{33}v_3 = \begin{bmatrix}0\\0\\0\end{bmatrix}.
\]
Il processo termina perché $h_{43}=0$.

\end{frame}
\begin{frame}{Risultato finale:}
\[
V_3 = [v_1\ v_2\ v_3] =
\begin{bmatrix}
1 & 0 & 0\\
0 & 1 & 0\\
0 & 0 & 1
\end{bmatrix},
\quad
\underline{H}_3 =
\begin{bmatrix}
2 & 1 & 0\\
1 & 2 & 1\\
0 & 1 & 2\\
0 & 0 & 0
\end{bmatrix}.
\]

\end{frame}


\end{document}

